\section{项目名称与意义}
\begin{frame}{项目背景与实用价值}
\begin{block}{项目名称：Web 应用自主测试智能体}
面向自然语言测试场景，自主操作浏览器、验证检查点，并生成可追溯的测试报告。
\end{block}
\begin{itemize}
\item 测试人员需要把业务描述转换成操作步骤，并逐项核对页面结果。
\item 页面结构变化会增加脚本定位与维护工作，长流程也容易遗漏检查点。
\item 本项目拟将场景输入、执行跟踪和证据查看集中到中文工作台，减少重复操作。
\item 面向课程指定应用开展实验，评估自然语言驱动测试的适用范围与失败原因。
\end{itemize}
\end{frame}

\begin{frame}{拟探索的改进点与验证方式}
\begin{description}
\item[证据约束的判定] 保留原始检查点，要求结论引用页面观察与截图，独立审查证据是否直接支持要求。
\item[有限预算的恢复] 在最多 35 次浏览器操作内，结合页面反馈调整定位与动作，限制重复尝试。
\item[可追溯的过程] 串联场景、动作、检查点与报告，使测试人员能复核结论和失败原因。
\end{description}
\begin{block}{实验验证计划}
在相同场景与操作预算下，对比是否启用独立证据审查、是否启用有限恢复，统计判定一致率、无法判定比例与操作数。
\end{block}
\small 上述内容为工程改进与实验假设，效果需通过真实场景验证。
\end{frame}
